\chapter{Computations as decisions with costs}
\label{sec:p17-metareasoning}

Metareasoning asks which computation to perform, how much resource to
give it, and when to act on the answer already available. A computation
earns decision value through the external action it changes. Improving
an evaluation number while choosing the same action is not by itself
an improvement in the external outcome.

The source is Russell and Wefald's complete 1989 conference paper,
\emph{Principles of Metareasoning}, pp.~400--411
\cite{RussellWefald1989KRMetareasoning}. Its later 1991 journal article
is a different presentation. The 1989 paper gives formulas for complete
computations, two approaches to partial computations, and practical
approximations. It explicitly leaves the general partial-computation
problem and some questions about imperfect self-evaluation unresolved.
The finite constructions below supply precise conditional solutions
and proofs without identifying those solutions with a resolution of
the source's unrestricted problem.

\section{Computations, intentions and resources}
\label{sec:p17-meta-actions}

\begin{hbdefinition}{def:p17-meta-source-notation}{The source's actions and evaluation snapshots}
Let \(A_1,\ldots,A_m\), \(m\geq1\), be the available external actions,
\(S_j\) a computational action, and \(S\) the computation history.
The source writes \([X,W]\) for the situation resulting from
action \(X\) in situation \(W\), and \([X]\) when the starting
situation is understood. Its \(U(W)\) is actual external utility;
\(\widehat U^S(W)\) is an estimate formed after history \(S\).
They need not coincide. The default intention is
\[
 \alpha(S)=A_{i^*},\qquad
 i^*=\operatorname*{arg\,max}_{i}\widehat U^S([A_i]),
\]
with a specified tie rule. A complete computation commits to an external
action. A partial computation leaves a further choice between acting
and computing. Write \(T\) for a future computation sequence that
ends in an external action \(\alpha_{S.S_j.T}\).

For a specified continuation policy that eventually acts, with
integrable terminal utilities, the source's actual continuation
value has the form
\[
 U([S_j])=\sum_T P(T)\,
                     U([\alpha_{S.S_j.T},[S_j.T]]).
\]
The probabilities describe that policy's possible complete traces,
including their computation outcomes. A policy chosen before those
outcomes are known can depend on outcomes as they arrive; it cannot
select retrospectively whichever realized trace was best.
If a policy may never act, its nontermination payoff needs a separate
definition; a sum over completed traces alone need not have total mass one.
The time and space used by a computation are resources. Their utility
cost must come from a declared delay model, resource price, or feasibility
constraint; a number of search nodes alone is not a utility scale
\cite[\S\S2--3, pp.~401--403]{RussellWefald1989KRMetareasoning}.
\end{hbdefinition}

\begin{hbdefinition}{def:p17-meta-program-model}{A model induced by actual program steps}
A machine configuration contains its program, input, current workspace,
search frontier or other retained results, evaluation vector, elapsed
time \(t\), and memory use. A feasible computation \(c\) is a
specified executable block on that configuration. Its transition
\(F_c\) produces a successor configuration, its running time and its
peak memory requirement. The transition may be deterministic; executing
a deterministic program need not give its controller advance knowledge
of the answer.

A metalevel information state \(z\) is the declared summary available
to the controller: for example the retained tree and its backed-up
estimates, computation type, input features, clock and remaining memory.
Specify a finite model of problem instances \(\Omega\), with
conditional weights \(\mu_z(\omega)\), and a summary map \(I\).
For a deterministic block define the predictive kernel
\[
 K(z'\mid z,c)=\sum_{\omega:
       I(F_c(\omega,z))=z'}\mu_z(\omega).
\]
Here \(F_c(\omega,z)\) includes the instance-specific configuration
consistent with summary \(z\); the model must specify it, or include
the additional configuration uncertainty in \(\omega\).
Any physical random bits used by the program may also be included in
this finite model. Require \(\mu_z(\omega)\geq0\) and
\(\sum_\omega\mu_z(\omega)=1\); every permitted block finishes
on the modeled support. Execution obtains no new external observation unless
one is expressly part of the action. Its output is a computed revision
of the workspace and decision estimates.

The conditional instance law is a performance model or a subjective
prediction over unresolved computations. It is not physical randomness
created by a deterministic calculation. Conditioning on complete input
and program under logical omniscience would make its answer fixed.
Thus an empirical law conditioned on a compressed summary must not be
claimed exact for every individual input. This is the distinction the
source makes between expensive unperformed computations and statistical
knowledge of similar completed computations
\cite[\S4.2, pp.~404--405]{RussellWefald1989KRMetareasoning}.
\index{metareasoning}\index{value of computation}\index{computational action}
\end{hbdefinition}

\begin{hblemma}{thm:p17-meta-kernel}{proved=annex:p17.meta.kernel}{Induced transitions and a sufficient information state}
The displayed kernel is a probability distribution on its reachable
summaries. For each observed successor of positive probability, the
conditional weights proportional to
\(\mu_z(\omega)\mathbf1_{\{I(F_c(\omega,z))=z'\}}\)
give the successor instance law, after applying any specified instance
transition. No such quotient is required at a zero-probability successor.
If, conditional on the entire observed history, the next-summary law,
available computations, their costs and the stopping payoff depend only
on \(z\), then \(z\) is sufficient for the finite recursion below.
If this condition fails, past computation information must be retained;
the Markov recursion on the coarser summary is unjustified.
\end{hblemma}

\section{Complete computations and consistent comparison}
\label{sec:p17-meta-complete}

\begin{hbdefinition}{def:p17-meta-complete-value}{New estimates and the same comparison baseline}
Fix the current intention \(\alpha\). After a complete computation
\(c\), let \(X=(X_1,\ldots,X_m)\) be its random vector of
new intrinsic action estimates, with specified joint distribution
\(p_c\). A named action can index its coordinate: \(X_{A_i}=X_i\).
Select \(\alpha_c\) by maximizing that new vector with
the declared tie rule. Assume, for this definition, that delaying all
external actions by this computation subtracts the same cost \(k_c\)
in utility units. Then the source's new-estimator comparison is
\[
 \widehat{\operatorname{VOC}}_1(c)
   =E_{p_c}\,[X_{\alpha_c}-X_\alpha]-k_c
   =E_{p_c}\,[\max_i X_i-X_\alpha]-k_c.
\]
Both candidates use the later intrinsic estimate. The old intention's
term is its estimated value if performed without the new delay, not
its delayed payoff. This is the separation behind source equations
(9), (11) and (15). With a density, replace the finite expectation by
the corresponding integrals, assuming integrability.

The common delay loss is an additional hypothesis. In the general
model use the actual action-dependent payoff at the successor clock.
For example, an opportunity to use action 1 can disappear while action 2
remains unchanged. No scalar time cost represents both losses, and
the preferred action may then change solely because of delay
\cite[\S\S4.1--4.3, pp.~404--406]{RussellWefald1989KRMetareasoning}.
\end{hbdefinition}

\begin{hbtheorem}{thm:p17-meta-baseline}{proved=annex:p17.meta.baseline}{The baseline correction and its calibrated special case}
The gross new-estimator benefit is nonnegative pointwise. If the
computation never changes the intention, its net score is \(-k_c\),
even if all estimates increase. Replacing the baseline by the earlier
estimate \(\widehat u_\alpha\) changes the expected score by
\[
                     E[X_\alpha]-\widehat u_\alpha.
\]
The two scores coincide exactly when this difference is zero.

If the new estimates are actual posterior expectations of fixed
integrable action utilities given the computed output, and the old
estimates are the corresponding prior expectations, then
\(E[X_i]=\widehat u_i\) for every action. When the old intention
maximizes those prior expectations, the gross benefit is exactly
the finite value-of-information improvement of
\cref{thm:p17-information-policy-value,thm:p17-information-value-bounds}.
This calibration condition supplies that connection; an arbitrary
heuristic evaluation function does not supply it.
\end{hbtheorem}

\begin{hbtheorem}{thm:p17-meta-subtree}{proved=annex:p17.meta.subtree}{The source's two subtree-independent formulas}
Assume common delay cost and that one computation changes only one
intrinsic action estimate.
If it changes a currently inferior action \(\beta\), while the
current best estimate remains \(a\), and its new estimate is \(X\),
the net score is
\[
                       E[(X-a)_+]-k_c,
                         \qquad x_+=\max(x,0).
\]
If there is at least one other action and the computation changes
the current intention's estimate to \(X\), while the
best other estimate remains \(b\), the net score is
\[
                              E[(b-X)_+]-k_c.
\]
For an integrable density \(p\), these expectations are respectively
\(\int_a^\infty (x-a)p(x)\,dx\) and
\(\int_{-\infty}^b(b-x)p(x)\,dx\), reproducing source
equations (17)--(18). Revising the preferred action downward can make
computation useful by revealing that it should be avoided.
With only one external action, the gross decision-change benefit is
zero and the complete-computation score is just the negative cost.

The one-affected-estimate premise is distinct from probabilistic
independence. If several estimates change, their joint law is needed.
Equal marginals alone do not determine the expected improvement.
For example, if two estimates are always jointly \((0,0)\) or
\((2,2)\), each with probability \(1/2\), the gross benefit for
either fixed old intention is zero. Making the two same marginals
independent would give gross benefit \(1/2\).
\end{hbtheorem}

\section{Grounding the partial-computation problem}
\label{sec:p17-meta-partial}

\begin{hbdefinition}{def:p17-meta-finite-control}{Feasible computational control with a terminal action}
Let \(Z\) be a nonempty finite set of sufficient summaries and
\(C(z)\) a finite set of permitted executable blocks. It may be empty.
The stopping option commits to the implemented intention \(\alpha(z)\),
with model expected external payoff \(g(z)\). Specifically, for
an external payoff model \(u(a,\omega,t)\), set
\[
             g(z)=\sum_\omega\mu_z(\omega)
                        u(\alpha(z),\omega,t(z)),
\]
where \(t(z)\) is the clock in that summary.
It does not secretly
optimize an uncomputed true-utility table. If several already available
external decisions can be chosen, their finite maximum may instead be
included as an explicitly computed stopping payoff.

The data are \(K(z'\mid z,c)\), expected direct resource cost
\(k(z,c)\), stopping payoff \(g(z)\), and a nonnegative integer
computation budget \(N\). A controller may stop at any time and must
stop after at most \(N\) blocks. Its return is its terminal external
payoff minus all accumulated direct computation costs. Clock-dependent
utility is already in \(g\) and the successor state; do not charge
the same delay again in \(k\). Under the separable time-cost model
one may instead keep intrinsic terminal payoff and charge its delay
loss once as part of \(k\).

Peak-space limits and deadlines remove a block from \(C(z)\) if
any positive-probability modeled outcome violates its declared resource
limit. An allowed block consumes its budget step and updates retained
workspace, clock and memory as specified by its program transition.
The metalevel's own evaluation cost must be included if incurred.
Fresh online evaluation of these options is itself a computation:
include its transition and cost, or apply the separately declared
controller invocation charge in \cref{thm:p17-meta-regress}.
The mathematical maxima below are model-based benchmark values;
their existence does not make an online evaluator free. A compiled
policy can instead execute a previously selected decision at its
specified lookup and branch cost.
All comparisons below are conditional on these supplied transition,
payoff and cost models. Their mathematical availability does not prove
that computing the comparison tables is cheap.
\end{hbdefinition}

\begin{hbtheorem}{thm:p17-meta-bellman}{proved=annex:p17.meta.bellman}{Exact finite computational control and stopping}
For every allowed randomized history policy define its expected
return by the finite products of its policy probabilities and \(K\).
The optimal return with at most \(n\) remaining blocks is
\[
 \begin{aligned}
 V_0(z)&=g(z),\\
 Q_n(z,c)&=-k(z,c)+\sum_{z'}K(z'\mid z,c)V_{n-1}(z'),\\
 V_n(z)&=\max\{g(z),\ \max_{c\in C(z)}Q_n(z,c)\}.
 \end{aligned}
\]
When \(C(z)\) is empty, use just \(g(z)\). A deterministic
policy choosing a maximizing stop or computation option is optimal
against all randomized history policies. Stop is optimal exactly when
every \(Q_n(z,c)\leq g(z)\); a strict reverse inequality justifies
continuation. Ties may be resolved in favor of stopping.
For finite rational tables this recursion and its selected policy are
computable exactly. It grounds the source's future-computation values
in the explicit budget boundary, using the finite-horizon method of
\cref{thm:p17-sequential-backward}.
\end{hbtheorem}

\begin{hbdefinition}{def:p17-meta-policy-model}{A prediction of the actual continuation policy}
One can predict a controller without assuming that it is optimal.
For remaining budget \(n>0\), let \(p_n(z)\) be its stopping
probability and \(d_n(c\mid z)\) its conditional choice distribution
if it continues. Then
\[
 W_0(z)=g(z),\qquad
 W_n(z)=p_n(z)g(z)+(1-p_n(z))
       \sum_c d_n(c\mid z)
       [-k(z,c)+\sum_{z'}K(z'\mid z,c)W_{n-1}(z')].
\]
If continuing is impossible, set \(p_n(z)=1\) and omit the
computation sum. A history-dependent version retains that history.
This is the finite probability-of-acting self-model of the source's
equations (21)--(22), with a stated stopping boundary. Expanding it
sums over complete computational traces; their final actions are
weighted by the probability of being selected. Estimated utilities
must be conditioned on that selection when it carries information.
\end{hbdefinition}

\begin{hbtheorem}{thm:p17-meta-policy-law}{proved=annex:p17.meta.policy-law}{Trace expansion and why selection changes an estimate}
The policy recursion equals expected terminal payoff minus accumulated
cost, summed over its complete traces, and \(W_n\leq V_n\).
For any random selected action \(A\) and actual terminal utility
\(U_A\) in this finite model,
\[
 E[U_A]=\sum_{a:\Pr(A=a)>0}\Pr(A=a)E[U_a\mid A=a].
\]
Replacing the conditional means by unconditional means is generally
false. If two action utilities are \((1,0)\) or \((0,1)\),
with equal probabilities, a controller selecting the larger has value
one. Averaging its selection probabilities against the unconditional
means instead gives \(1/2\).
\end{hbtheorem}

\section{Approximations, overhead and regress}
\label{sec:p17-meta-approximation}

\begin{hbmethod}{def:p17-meta-feasible-controller}{A compiled computation policy}
Before deployment, record completed episodes by computation type and
declared input/workspace features. Estimate their joint changes to the
evaluation vector, their running time, peak space and the subsequent
external choice. Use the same completed-episode evaluation of the new
and old intention. Store a transition model or a direct score table
indexed by those features. Online, use a fixed finite lookup, a resource
feasibility check, and a stated stopping rule. Price those operations.

A one-step controller uses \(V_0=g\) in every computation comparison:
it evaluates a step as though action followed immediately. The source
calls this its single-step assumption. Restricting candidates to individual
node expansions is its separate meta-greedy assumption. They can make
control inexpensive, but discard complementarity between partial steps.
The source's adaptive proposal backs up completed episode evaluations
and revises their distributions with experience. Its convergence is a
suggestion, not an established theorem; changes to the controller change
which episodes it generates \cite[\S\S4.2, 4.5--5]{RussellWefald1989KRMetareasoning}.
In particular, the source's proposed sample-label construction provisionally
assumes that the current stop/continue decisions are correct while backing
up completed computations. If those decisions are poor, the labels evaluate
that policy's consequences; they do not automatically evaluate an optimal
continuation. The grounded \(W_n\) construction states exactly which
policy is being evaluated. Treating its values as \(V_n\) requires
the additional optimization or a justified approximation bound.
\end{hbmethod}

\begin{hbtheorem}{thm:p17-meta-approximate}{proved=annex:p17.meta.approximate}{What an approximation certificate actually guarantees}
At every state with \(n\) blocks left, suppose approximate scores
\(\widetilde Q_n(z,a)\) for every option, including stop, differ
from their exact optimal-continuation scores by at most \(\epsilon_n\).
The policy selecting a maximum approximate score and applying the same
rule thereafter has loss at most
\[
                         2\sum_{j=1}^n\epsilon_j.
\]
An approximate greedy score advantage greater than twice the applicable
error bound preserves the corresponding exact comparison.
For one-step stopping, if all exact net values are nonpositive and
their one-step approximations have error at most \(\epsilon\),
any accepted score exceeding \(\epsilon\) is inconsistent with
that error guarantee. Conversely an estimated one-step net value
greater than \(\epsilon\) certifies positive exact one-step value.
It does not certify globally optimal allocation of further computation.
\end{hbtheorem}

\begin{hbtheorem}{thm:p17-meta-regress}{proved=annex:p17.meta.regress}{Ungrounded continuation and the cost of choosing a computation}
The equation \(w=\max(g,w)\), representing a zero-cost unchanged
continuation with no terminal boundary, has every \(w\geq g\)
as a solution. Its positive-cost version
\(w=\max(g,w-c)\), \(c>0\), has the sole finite solution \(g\).
Thus a recursive valuation formula alone need not determine a value.

If using a computation controller costs an action-independent amount
\(r\geq0\) before its chosen computation or stop option, its value
is its policy return minus \(r\). Compared with taking the current
intention directly, the controller is worthwhile exactly when its
policy improvement exceeds \(r\), with equality giving a tie.
For example, a complete computation with gross benefit one and direct
cost \(3/4\) has net benefit \(1/4\). Paying \(3/10\) to
select it changes the net benefit to \(-1/20\).
Further deliberation about that controller needs its own priced action
and stopping boundary; assuming a free optimizing metalevel does not
solve the evaluation regress.
\end{hbtheorem}

\section{A calculation with an actual two-block computation}
\label{sec:p17-meta-computation-example}

\begin{hbdefinition}{def:p17-meta-two-block-program}{A retained intermediate result and a terminal evaluation}
Consider a fixed computational architecture with two input bits already
in its memory. Its controller's initial summary does not include their
values. The performance model assigns equal weights to the four input
pairs \((p,q)\in\{0,1\}^2\). These are problem instances, not
fresh external measurements. External action \(A\) has payoff zero.
The evaluation program for \(B\) returns utility four when the bits
agree and minus two otherwise. Its prior estimate is therefore one;
the initial intention is \(B\).

Block \(c_1\) loads \(p\) into retained workspace at cost
\(2/5\). Block \(c_2\), available afterwards, loads \(q\),
computes the exclusive-or of the two stored bits, and returns the
exact evaluation of \(B\), at cost \(2/5\). These are the two
available computational blocks; each has a declared one-unit duration,
and both fit the memory limit. There is no further delay loss in this
example. The terminal choice takes \(B\) after an equal-bit result
and \(A\) after an unequal-bit result.

This is a model of a specified instruction and control interface.
It makes no claim that exclusive-or is intrinsically expensive or that
the supplied feature law is exact for every other machine. The controller
cannot inspect an unreported input component without executing a permitted
priced block. Other instruction sets, summaries or costs give other
metareasoning problems.
\end{hbdefinition}

\begin{hbtheorem}{thm:p17-meta-two-block}{proved=annex:p17.meta.two-block}{Stopping, complementarity and a feasible compiled policy}
With one block available from the initial state, stopping has value one
and performing \(c_1\) then stopping has value \(3/5\).
The one-step controller stops. With two blocks available, the exact
value is \(6/5\): perform \(c_1\), perform \(c_2\), then act
according to the returned evaluation. At either possible first-block
output the value of the second block followed by action is \(8/5\),
exceeding the current stopping value one.

Thus both individual first-block outputs leave the action estimate
unchanged, while the two-block computation has net benefit \(1/5\).
A compiled policy implementing these two blocks and the terminal choice
with an additional total control cost \(1/10\) has return \(11/10\)
and still improves on immediate action. If only one block fits a deadline,
the two-block policy is infeasible and stopping is optimal.
The result is an explicit instance of the source's meta-greedy barrier:
one-step nonimprovement does not imply that a sequence is worthless.
\end{hbtheorem}

\section{Applications and the limits of the reported evidence}
\label{sec:p17-meta-applications}

In the source's search examples, a computation expands frontier nodes
and backs their new values up through a retained search tree. Predicting
the new leaf values is empirical; transmitting them through the chosen
backup procedure is analytic. Shared graph nodes or a probability revision
that affects many alternatives can invalidate the one-subtree premise.
The paper distinguishes three qualitative cases: estimates so separated
that a change is unlikely; near symmetry with little attainable payoff
difference; and overlapping uncertain alternatives with useful expected
change. Distance between two estimated means alone does not settle any
of these cases; the predictive law and computation cost matter.

The 1989 applications use the single-step approximation. DTA* compares
the external path cost with a declared search-node/motion cost ratio.
The paper's Figure~4 shows eight-puzzle and 20-polygon path-planning comparisons
against A* and fixed-horizon RTA*. The figure is reported performance
over varying cost ratios, not a universal optimality theorem. Details
of the algorithms and search-theoretic proofs are referred by the source
to its separate 1988 technical report.
At the stated admissible-heuristic setting, the conference paper also
attributes to that report a restriction of useful expansion to the
current best move's subtree, a computable benefit for a set of expansions,
a best-first policy given only the heuristic information, and a real-time
stopping criterion. These attributed results retain their source-model
conditions; the reported comparison does not prove that A* minimizes
total physical-action-plus-computation cost in every metalevel model.

MGSS* additionally uses individual-node meta-greedy comparisons.
It can stop when no one-node expansion changes the preferred move,
although a coordinated set would change it. The paper's Table~1 reports one
32-game Othello tournament against depth-three alpha--beta search:
16 random starting positions, each played with reversed colors.
MGSS* won 22 games and alpha--beta won 10. The table's time entries
are 1095 versus 1261 seconds and its node entries are 190077 versus
279697; it does not label those two rows with a per-game or whole-tournament
denominator. The paired game setup and reported entries do not establish statistical
independence of 32 games, a confidence interval, reproducibility on
unreported positions, or comparison with every search budget. The paper
gives no instance records or machine specification sufficient here to
reproduce the tournament \cite[\S5, pp.~408--410]{RussellWefald1989KRMetareasoning}.

The source's future-work discussion proposes retaining information
between deliberations, recursive resource allocation to reduce memory,
and a broader decision-theoretic architecture. The theorem-proving
proposal needs a usable intermediate recommendation and a utility for
partial progress. Estimated success or utility cannot be substituted
for an actual proof. These are research directions in that presentation.
Likewise its probabilistic self-model and adaptive method do not settle
every philosophical question about an agent trusting its own imperfect
estimates. A finite priced model, a grounded continuation rule, and a
validated prediction law are the inputs required to apply the proved
control results of this chapter
\cite[\S\S4.5, 6--7, pp.~407--411]{RussellWefald1989KRMetareasoning}.
